\section*{Chapter \ref{ch:nw:capturing-and-monitoring-the-network} --- Capturing and Monitoring the Network}

\begin{solution}{exo:nw:capturing-and-monitoring-the-network:1}
No: the nanoseconds field must be below $10^9$. A lenient reader would carry the second and report 34\,201~s and 500\,000\,321~ns,
09:30:01.500\,000\,321; a strict one should reject the record, because the writer is broken.
\end{solution}

\begin{solution}{exo:nw:capturing-and-monitoring-the-network:2}
$16\,771 - 12\,345 = \qty{4\,426}{\nano\second}$, from the first copy (line~A). With line~A's copy lost the strategy saw the packet on line~B
\qty{3}{\micro\second} later and its order would have left about \qty{3}{\micro\second} later: measured from the B copy the wire-to-wire latency is
the same, and the \qty{3}{\micro\second} belong to the lost line.
\end{solution}

\begin{solution}{exo:nw:capturing-and-monitoring-the-network:3}
Each direction carries $4 \times 10^9 / (220 \times 8) \approx 2.27$ million frames a second (200 bytes plus 20 of preamble and gap on the
wire); each is stored as $200 + 16 + 16 = 232$ bytes: about 527~MB a second per direction, 1.05~GB a second for the link.
\end{solution}

\begin{solution}{exo:nw:capturing-and-monitoring-the-network:4}
The rule assigns an order to the latest trigger before it. It is wrong whenever another trigger arrives between the true trigger and the
order, which happens with a probability that grows with the trigger rate times the firm's latency: frequent triggers, or a slow firm, put more
triggers in the window.
\end{solution}

\begin{solution}{exo:nw:capturing-and-monitoring-the-network:5}
$99\,055 \times 0.002^2 \approx 0.4$: seeing none is what independence predicts. A packet lost on both lines is a common cause.
\end{solution}

\begin{solution}{exo:nw:capturing-and-monitoring-the-network:6}
The median: the card's two passes (1.6~\unit{\micro\second} together, a better card or kernel bypass on a tuned host) and the ring
(a busy-polling consumer on an isolated core). The 99th percentile: the ring, whose tail comes from scheduling; isolating its cores costs
cores and tuning (One Quant Book~13, chapter~13), not money.
\end{solution}

\begin{solution}{exo:nw:capturing-and-monitoring-the-network:7}
The naive rule's error falls from 47\% to 25\%: with a four times faster firm, fewer triggers arrive between a trigger and its order. It does
not vanish, because \omterm{def:nw:capturing-and-monitoring-the-network:trigger}{trigger packets} still come in bursts closer together than the firm's latency.
\end{solution}

\begin{solution}{exo:nw:capturing-and-monitoring-the-network:8}
The card's timestamps miss everything outside the card: the cage's fibre and devices on both sides, which the tap sees. They also rely on the
card's clock and on how the card defines its receive and transmit instants. It is a good in-host measurement; it is not the firm's
wire-to-wire latency at the handoffs.
\end{solution}

\begin{solution}{pb:nw:capturing-and-monitoring-the-network:1}
\begin{enumerate}
\item $437 + 756 + 158 + 97 + 47 + 33 = \qty{1\,528}{\nano\second}$.
\item $1\,528 + 263 + 295 + 798 + 801 = \qty{3\,685}{\nano\second}$.
\item The stages are skewed to the right: the median of a sum of skewed variables lies above the sum of their medians (here
  \qty{4\,386}{\nano\second} in the budget).
\item $4.4 - 1.5 = \qty{2.9}{\micro\second}$: two thirds of the median.
\item $31.9 / 38.5 = 83\%$.
\item A high percentile of the sum is dominated by the stage whose tail is heaviest; the other stages' tails rarely coincide with it.
\item A 99.9th percentile of \qty{3.5}{\milli\second}: the consumer descheduled on shared cores.
\item Isolated cores and a busy-polling consumer, which remove the scheduler from the ring.
\item The tap sees the frame where the firm receives it from the venue; the card's timestamp comes after the cage's network, so it misses
  part of the path, and the card's clock is one more thing to trust.
\item Understated: the naive rule pairs slow orders with later packets, and gives a 99th percentile of \qty{27}{\micro\second} instead of 38.
\item The client order identifier, or a field the venue returns unchanged.
\item Wire-to-wire percentiles computed from the capture every minute, compared with the baseline.
\item \textbf{Named result.} Wire-to-wire median \qty{4.4}{\micro\second} and 99th percentile \qty{38}{\micro\second}; \qty{2.9}{\micro\second} of the
  median lies outside the program; the ring between handler and strategy owns the tail.
\item By about \qty{0.9}{\micro\second} (the budget's median falls from 4.39 to 3.47).
\item Decode (p99 \qty{2.8}{\micro\second}), with the card's two passes close behind at \qty{2}{\micro\second} each.
\item Hardware timestamps at the card's port on both sides, compared with the taps', or a capture on each side of the server.
\item It is what the firm actually sent and received, and when, independent of any program's logs: the record a venue, a client or a
  regulator asks for.
\item About 38~MB a second of capture, 138~GB an hour.
\item A month of summaries (latency percentiles, gaps, bursts, alarms) and of the order handoff's frames; a day of every market-data frame,
  unless an incident asks for more.
\item Everything outside the program, timed by a clock the program cannot move.
\end{enumerate}
\end{solution}

\begin{solution}{iq:nw:capturing-and-monitoring-the-network:1}
Tap the market-data handoff and the order handoff, timestamp both on the wire against one disciplined clock, pair each order with its \omterm{def:nw:capturing-and-monitoring-the-network:trigger}{trigger
packet}, and subtract. A program's logs miss the network, the card and the kernel on both sides, and use the host's clock.
\iqlookfor{taps, one clock, and pairing.}
\end{solution}

\begin{solution}{iq:nw:capturing-and-monitoring-the-network:2}
On the wire, at the first byte of the frame, by the tapping device, carried with the copy (a trailer). Anything later adds the aggregation
network's queueing and the storage server's clock to every measurement.
\iqlookfor{timestamp at the tap, not at storage.}
\end{solution}

\begin{solution}{iq:nw:capturing-and-monitoring-the-network:3}
Make the strategy carry the trigger's sequence number (or a hash) in the order, in a field the capture can read. Heuristics such as the
latest trigger before the order fail in bursts, which is when latency matters.
\iqlookfor{instrumenting the order rather than guessing.}
\end{solution}

\begin{solution}{iq:nw:capturing-and-monitoring-the-network:4}
Egress discards on ports that should never be full, gaps per line and on both lines, snapshot recoveries, clock lock and offsets, and
wire-to-wire percentiles against a baseline: each points at a place on the path.
\iqlookfor{alarms tied to mechanisms, not a dashboard of everything.}
\end{solution}

\begin{solution}{iq:nw:capturing-and-monitoring-the-network:5}
Decompose it: compare the capture's stage timestamps (taps, card timestamps if kept) with the budget; check whether the market got busier
(bursts, message rates), whether a clock moved (\omterm{def:nw:time-synchronisation:holdover}{holdover}, a new asymmetry), whether a host changed (kernel, firmware, a new process on a
shared core), and whether the venue changed its feed. Look at when it started to the second.
\iqlookfor{decomposition by stage and a list of non-deployment causes.}
\end{solution}

\begin{solution}{iq:nw:capturing-and-monitoring-the-network:6}
Gap and loss analysis, \omterm{def:nw:networking-for-trading:burst}{microburst} detection, rebuilding the book to check the handler, reconstructing what the firm sent during an incident,
verifying venue acknowledgement times, and evidence for disputes and regulators.
\iqlookfor{the capture as evidence and as a second, independent implementation.}
\end{solution}
